\chapter{Conclusion}

\begin{singlespace}
This report presented a lightweight anonymous authentication scheme for multihop
IoT networks that extends DAuth~\cite{das2026comsnets} with two additions:
per-session pseudonym rotation to hide $ID_D$ from the open channel, and a dual-state
mechanism to recover from desynchronisation without re-enrolment.

Formal verification in ProVerif confirms all 10 security queries under the Dolev--Yao
model, covering authentication correspondence, session key secrecy, forward secrecy,
anonymity, and offline guessing resistance.

COOJA\slash Contiki-NG simulations on 100-node networks show per-device energy
reductions of \textbf{42.8\,\%} over LAAKA and \textbf{32.4\,\%} over Zhou et al.
The dual-state mechanism is particularly significant in the desynchronised case:
DAuth incurs a \textbf{118.3\,\%} energy penalty per recovery round
(134.8\,mJ vs.\ 61.8\,mJ for a normal round) because a single lost packet forces
full re-enrolment, whereas the proposed scheme's recovery round costs 87.2\,mJ
(close to its normal round of 98.6\,mJ), saving \textbf{35.3\,\%} in both
energy and CPU time over DAuth.

Hardware experiments on Raspberry Pi~4B confirm per-round energy reductions of
\textbf{52.5\,\%} over LAAKA and \textbf{14.2\,\%} over Zhou et al., with an
88\,\% overhead above DAuth attributable solely to the additional hash operations
for pseudonym rotation.

Among the four compared schemes, the proposed scheme is the only one to simultaneously
provide device anonymity, desynchronisation recovery, and decoupled distributed
authentication at overhead suitable for resource-constrained IoT devices.
\end{singlespace}
